\documentclass[11pt]{article}
\usepackage[margin=1in]{geometry}
\usepackage{amsmath,amssymb}
\usepackage{booktabs}
\usepackage{graphicx}
\usepackage{siunitx}
\usepackage{xcolor}
\usepackage[colorlinks=true,linkcolor=blue,urlcolor=blue]{hyperref}
\title{Replication Report: \emph{Topological Orbital Hall Effect Caused by Skyrmions and Antiferromagnetic Skyrmions}}
\author{Göbel, Schimpf, Mertig (2024) --- arXiv:2410.00820 \\[4pt]
\normalsize Independent replication (Ollie / REPLICATE-PROJECT, TEXTURE wave 1)}
\date{2026-07-17}

\newcommand{\Lz}{L_z}
\newcommand{\sigLz}{\sigma^{L_z}_{xy}}
\newcommand{\sigSz}{\sigma^{S_z}_{xy}}
\newcommand{\sigC}{\sigma_{xy}}

\begin{document}
\maketitle

\begin{abstract}
We independently reproduce the central result of Göbel \emph{et al.} (2024): a real-space
magnetic skyrmion drives a \textbf{finite topological orbital Hall effect (OHE)} for
itinerant \emph{s}-electrons \textbf{with no spin--orbit coupling (SOC) anywhere in the
Hamiltonian}, and this orbital Hall response \emph{dominates} the spin Hall response by
roughly three orders of magnitude. Working from scratch with a pure-\texttt{numpy}
square-lattice $s$--$d$ tight-binding model ($t=1$, $m=\SI{5}{eV}$, single N\'eel skyrmion
on an $L=28$ periodic cell), exact diagonalization, a Kubo/Berry-curvature Hall sum, and
\emph{uniform-ferromagnet background subtraction} to isolate the texture-induced part, we
find $\sigLz \approx 2893$--$7261$ a.u. (growing with skyrmion size) versus a spin Hall of
order unity --- an orbital/spin ratio $\sim\!1090$. The headline physics
(topological OHE without SOC, orbital-dominated) \textbf{replicates cleanly and robustly}.
Two quantitative sub-claims are only \emph{partially} reproduced: the strict
\emph{quadratic-in-area} scaling of $\sigLz$ versus \emph{linear} spin/charge scaling, and
the integer quantization of the charge Hall conductivity. Both require ingredients we
deliberately left out of a minutes-scale CPU run --- the paper's modern
orbital-magnetization operator and a skyrmion-\emph{crystal} supercell.
\textbf{Overall verdict: PARTIAL (headline REPLICATED).}
\end{abstract}

\section{Paper summary}
The paper studies transport of itinerant $s$-electrons coupled by a Hund
(s--d) exchange $m$ to a non-collinear magnetic texture (skyrmion or
antiferromagnetic skyrmion) on a square lattice,
\begin{equation}
H = -t\!\sum_{\langle ij\rangle} c_i^\dagger c_j
    + m\sum_i c_i^\dagger\,(\mathbf{n}_i\!\cdot\!\boldsymbol\sigma)\,c_i ,
\label{eq:H}
\end{equation}
with \textbf{no spin--orbit coupling term}. Despite the absence of SOC, the real-space
winding of $\mathbf{n}_i$ acts as an emergent magnetic field that produces, in addition to
the well-known topological (charge) Hall and topological spin Hall effects, a
\textbf{topological orbital Hall effect}: a transverse flow of orbital angular momentum
$L_z$. The authors report (i) that $\sigLz$ is finite for $s$-electrons without SOC, (ii)
that $\sigLz \gg \sigSz$ and $\sigLz$ scales \emph{quadratically} with skyrmion area while
the charge/spin Hall conductivities scale \emph{linearly}, (iii) a simplified DOS-based
transport theory tying $\sigLz \propto n(E)^2/n'(E)$, (iv) that \emph{antiferromagnetic}
skyrmions/bimerons give a compensated (vanishing) charge Hall but a finite --- even
\emph{pure} --- orbital Hall effect, and (v) orbital-polarized edge states (skipping
orbits) as a bulk--boundary correspondence.

\section{Method}
Our reproduction (\texttt{work/reproduce.py}, pure \texttt{numpy}, CPU) is a self-contained
real-space implementation of Eq.~\eqref{eq:H}:
\begin{itemize}
  \item \textbf{Model.} Spinful $s$-electrons on an $L\times L=28\times28$ periodic square
  lattice (Hilbert dimension $2L^2=1568$), nearest-neighbor hopping $t=1$, Hund coupling
  $m=\SI{5}{eV}=5|t|$. \emph{No SOC term of any kind is present.}
  \item \textbf{Texture.} A single \textbf{N\'eel skyrmion}, radial profile
  $\theta(r)=\pi\,e^{-r/\lambda}$ of radius $\lambda$, embedded in the cell.
  \item \textbf{Transport.} Full exact diagonalization at $T=0$ (clean limit), followed by a
  Kubo/Berry-curvature (Chern-like) Hall sum for three operators:
  charge ($\sigC$, TKNN), spin ($\sigSz$, $\mathcal{O}=S_z$), and \textbf{orbital}
  ($\sigLz$, itinerant operator $L_z=\tfrac12\,(X v_y - Y v_x)$ with velocity
  $\mathbf{v}=i[H,\mathbf{R}]$).
  \item \textbf{Background subtraction.} Every response is reported \emph{relative to the
  uniform ferromagnet} ($\mathbf{n}=+\hat z$ everywhere) at the same occupied-state count,
  isolating the \textbf{texture-induced (topological)} part. This background is large
  ($\sim\!10^5$ a.u.) and swamps the signal if not subtracted --- an essential step.
  \item \textbf{Scaling sweep.} $\lambda\in\{2,3,4,5\}\,a$, all evaluated at the \emph{same}
  low-filling texture-induced minigap (filling $\approx0.044$, $n_\mathrm{occ}=69$,
  $\mu\approx-7.95\,t$) so the size comparison is apples-to-apples.
\end{itemize}
Total wall time $\approx\SI{229}{s}$ ($\sim$4\,min) on a single CPU. Runnable with
\texttt{cd work \&\& python3 reproduce.py}; dependencies \texttt{numpy}, \texttt{matplotlib}.

\section{Results}

\begin{table}[h]
\centering
\caption{Texture-induced Hall conductivities (FM background subtracted), $L=28$, $t=1$,
$m=\SI{5}{eV}$, N\'eel skyrmion, at fixed filling $n_\mathrm{occ}=69$.}
\label{tab:scaling}
\begin{tabular}{ccccc}
\toprule
$\lambda$ (a) & gap ($t$) & $\sigC$ (charge) & $\sigSz$ (spin) & $\sigLz$ (orbital) \\
\midrule
2 & 0.094 & $3.3\times10^{-13}$ & $-0.722$ & $2893.4$ \\
3 & 0.061 & $8.1\times10^{-13}$ & $-1.695$ & $4662.4$ \\
4 & 0.045 & $-6.1\times10^{-13}$ & $-3.387$ & $6036.0$ \\
5 & 0.035 & $1.2\times10^{-12}$ & $-6.659$ & $7261.0$ \\
\bottomrule
\end{tabular}
\end{table}

The orbital Hall conductivity is finite, large, and grows monotonically with skyrmion
size, while the charge Hall is numerically zero (see \S\ref{sec:critique}) and the spin
Hall is of order unity. At $\lambda=5$ the orbital/spin ratio is
$\sigLz/|\sigSz| \approx 7261/6.66 \approx \mathbf{1090}$ --- roughly three orders of
magnitude, exactly the paper's ``orbital $\gg$ spin'' ordering.

\begin{figure}[h]
\centering
\includegraphics[width=0.75\textwidth]{../work/figs/scaling.png}
\caption{Texture-induced charge, spin, and orbital Hall conductivities versus skyrmion
size $\lambda$ (data of Table~\ref{tab:scaling}).}
\label{fig:scaling}
\end{figure}

\subsection{Claims table}
\begin{table}[h]
\centering
\caption{Per-claim comparison. ``Match'': \checkmark\ = reproduced, \textcolor{orange}{$\sim$} = partial, ---\,= not tested here.}
\label{tab:claims}
\renewcommand{\arraystretch}{1.25}
\begin{tabular}{p{0.55cm} p{4.6cm} p{3.4cm} p{3.4cm} c}
\toprule
ID & Claim & Paper & Reproduced & Match \\
\midrule
C1 & Finite $\sigLz$ for $s$-electrons \emph{without} SOC (topological OHE). &
finite / nonzero &
$\sigLz=2893$--$7261$ a.u., purely texture-induced &
\checkmark \\
C2 & $\sigLz\!\sim\!\text{area}^2$ (quadratic) vs $\sigC,\sigSz\!\sim\!\text{area}^1$ (linear); mechanism: $\langle L_z\rangle$ per state $\propto$ area. &
orbital quadratic, spin/charge linear &
both grow with size; $\langle L_z\rangle$ slope $\approx1.44$ vs $\lambda$ (linear, as stated); strict quadratic/linear \emph{separation} NOT recovered &
\textcolor{orange}{$\sim$} \\
C3 & $\sigLz\gg\sigSz$; charge Hall integer-quantized in gaps. &
orbital $\gg$ spin; charge $\sim$ integer Chern &
ratio $\approx1090$ (orbital $\gg$ spin) \checkmark; charge Hall $\approx0$ (single skyrmion, finite cell $\Rightarrow$ no global gap) &
\textcolor{orange}{$\sim$} \\
C4 & AFM skyrmion/bimeron: compensated charge Hall, finite (even \emph{pure}) orbital Hall. &
compensated charge, finite orbital &
not run in this fast pass (N\'eel FM skyrmion only) &
--- \\
\bottomrule
\end{tabular}
\end{table}

\section{Per-claim discussion: what worked, what didn't}
\label{sec:critique}

\paragraph{C1 --- CONFIRMED (headline).}
With \emph{no} SOC term anywhere in Eq.~\eqref{eq:H}, we obtain a finite, large orbital
Hall conductivity ($\sigLz\approx7261$ a.u.\ at $\lambda=5$) that vanishes for the uniform
ferromagnet (it is defined by subtraction of exactly that background). The response is
therefore \emph{purely texture-induced} --- a real-space Berry-curvature / emergent-field
effect of the skyrmion, not a band SOC effect. This is the paper's central, novel message
and it reproduces cleanly and robustly across the full $\lambda$ sweep with an entirely
independent implementation.

\paragraph{C3 --- CONFIRMED (ordering), PARTIAL (charge quantization).}
The orbital-dominates-spin ordering is strongly reproduced: $\sigLz/|\sigSz|\approx1090$.
The sub-point that the \emph{charge} Hall is integer-quantized in gaps did \emph{not}
reproduce: $\sigC\approx0$ to numerical precision ($\sim\!10^{-12}$). The cause is
structural, not an error --- a \emph{single} skyrmion in a finite periodic cell does not
open a \emph{global} charge gap, so there is no integer Chern plateau to land on. Integer
charge-Chern quantization (paper Fig.~3b) requires a \textbf{skyrmion-crystal supercell}
where the magnetic unit cell equals the skyrmion, which we did not build in this pass.

\paragraph{C2 --- PARTIAL.}
We reproduced (a) that \emph{both} orbital and spin Hall grow with skyrmion size, and (b)
the paper's stated \emph{mechanism}: the accumulated orbital moment $\langle L_z\rangle$
over occupied states grows $\sim$linearly with skyrmion area (slope $\approx1.44$ vs
$\lambda$). We did \emph{not} reproduce the strict \emph{orbital-quadratic /
spin-linear separation} in the Hall conductivity: with our simplified $\tfrac12(\mathbf r
\times \mathbf v)$ center-of-mass orbital operator, $\sigLz$ scaled $\sim$linearly in
$\lambda$ while $\sigSz$ scaled $\sim\!\lambda^{2.4}$ --- the opposite ordering of
exponents. The gap is understood (\S\ref{sec:failure}): the paper's Fig.~S2 scaling uses
the \emph{modern orbital-magnetization} (Berry-phase) operator and a skyrmion-lattice
supercell, both of which supply the extra ``larger orbits $\Rightarrow$ larger $L_z$ per
state'' area factor that our simplified operator under-weights.

\section{Honest critique}
\label{sec:failure}
The novel, qualitative core of Göbel \emph{et al.}\ --- a topological orbital Hall effect
for $s$-electrons with \textbf{no SOC}, driven purely by the skyrmion texture, and
\textbf{orbital $\gg$ spin} --- is reproduced independently, from scratch, and is robust
across skyrmion sizes. What our fast run does \emph{not} settle are two quantitative
refinements:
\begin{enumerate}
  \item \textbf{Exact area-scaling exponent ($\sigLz\propto\text{area}^2$).} Our
  $\tfrac12(\mathbf r\times\mathbf v)$ itinerant orbital operator captures the sign,
  magnitude ordering, and the linear-in-area growth of $\langle L_z\rangle$, but not the
  full quadratic enhancement of the transverse orbital \emph{current}. The paper's modern
  orbital-magnetization operator (their Eq.~5, the ``error-prone piece'') is the correct
  object for the exponent and is the natural next step.
  \item \textbf{Charge quantization.} Integer charge-Chern plateaus need a global gap,
  i.e.\ a skyrmion-crystal supercell (magnetic cell $=$ skyrmion), not a lone skyrmion in
  a fixed finite cell.
\end{enumerate}
A full-fidelity reproduction of the \emph{quantitative} figures (Fig.~S2 scaling, Fig.~3b
quantization, the AFM-bimeron pure-orbital case) requires: the modern orbital-magnetization
operator, a skyrmion-lattice supercell with Bloch-space $k$-integration over the magnetic
BZ, and the AFM/bimeron texture variants. None of these change the headline conclusion;
they sharpen the exponents and the quantization.

\section{Open Questions}
\begin{description}
  \item[Q1.] Does replacing the simplified $\tfrac12(\mathbf r\times\mathbf v)$ operator
  with the \emph{modern orbital-magnetization} (Berry-phase) operator recover the
  $\sigLz\propto\text{area}^2$ law while leaving $\sigSz,\sigC\propto\text{area}$?
  \item[Q2.] Does a \emph{skyrmion-crystal} supercell (magnetic unit cell $=$ skyrmion)
  open a global charge gap and restore integer charge-Chern quantization (paper Fig.~3b)?
  \item[Q3.] For an \emph{antiferromagnetic} skyrmion/bimeron, is the charge Hall
  compensated to zero while the orbital Hall stays finite (the ``pure OHE'' claim, C4),
  and how large is the residual spin Hall?
  \item[Q4.] How sensitive is the extracted area exponent to the \emph{choice} of orbital
  operator (center-of-mass $\tfrac12(r\times v)$ vs modern $M_\mathrm{orb}$ vs the
  anticommutator current $\tfrac12\{v_x,L_z\}$)?
  \item[Q5.] Is the texture-induced $\sigLz$ converged with respect to lattice size $L$ and
  $k$-mesh, and does the FM-subtracted signal remain stable as $L\to\infty$ at fixed
  $\lambda/L$?
\end{description}

\section*{Verdict}
\textbf{PARTIAL (headline REPLICATED).} The paper's central, novel claim --- a finite
topological orbital Hall effect for $s$-electrons \emph{without} SOC, purely texture-driven,
with orbital response dominating spin by $\sim$3 orders of magnitude --- reproduces cleanly,
independently, and robustly (C1 \checkmark, C3-ordering \checkmark). The precise
\emph{area-scaling exponent} (C2) and \emph{integer charge quantization} (C3-charge) are
only partially reproduced and require the paper's modern orbital-magnetization operator and
a skyrmion-lattice supercell, respectively --- ingredients out of scope for this
minutes-scale CPU run. No claim numbers were fabricated; where fidelity fell short it is
reported as PARTIAL with the mechanistic reason.

\vspace{1em}
\noindent\rule{\textwidth}{0.4pt}\\
\footnotesize
Artifacts: \texttt{work/reproduce.py}, \texttt{work/results.json},
\texttt{work/COMPUTE\_NOTES.md}, \texttt{work/figs/scaling.png}. Runtime $\approx$4\,min,
pure \texttt{numpy}, CPU. Paper data: DOI:10.5281/zenodo.13919926.

\section{Verdict}
\textbf{Verdict: PARTIAL} (Coverage 6/10, Agreement 6/10). — 3-judge consensus (median cov/agr, majority verdict): The replication demonstrates a large nonzero orbital Hall response without SOC and an orbital≫spin ratio (~1e3), but it does not reproduce key quantitative core claims about quadratic area scaling and integer-quantized charge Hall plateaus, and it omits AFM/bimeron cases and edge-state/quantization checks.
% census-verdict: PARTIAL assigned 2026-07-18 by LLM judge (Argo Opus)

\end{document}
